\chapter{Thiết kế hệ thống}
\label{chap:thiet-ke}

\section{Kiến trúc tổng thể}
BrewLite tổ chức mã nguồn trong một repository, với frontend tại
\sourcefile{apps/web/} và backend tại \sourcefile{apps/api/}. Backend là
modular monolith: các module được triển khai chung trong một ứng dụng
NestJS, sử dụng PostgreSQL làm nguồn dữ liệu nghiệp vụ.
\begin{figure}[H]
\centering
\begin{tikzpicture}[
    box/.style={draw,rounded corners=2pt,align=center,text width=11cm,
        minimum height=1.05cm,fill=gray!5,font=\small},
    flow/.style={-{Latex[length=2.5mm]},thick},node distance=1cm]
\node[box] (browser) {\textbf{Trình duyệt khách hàng}\\
Next.js / React -- Menu, tùy chọn, giỏ, phiên, checkout, lịch sử};
\node[box,below=of browser] (api) {\textbf{NestJS REST API -- /api/v1}\\
Auth / Catalog / Orders / Inventory / Payments / Promotions};
\node[box,below=of api] (db) {\textbf{PostgreSQL qua Prisma}\\
Migration -- transaction -- ràng buộc -- snapshot -- reservation};
\draw[flow] (browser) -- node[right,font=\footnotesize] {HTTP / JSON / JWT} (api);
\draw[flow] (api) -- node[right,font=\footnotesize] {Đọc / ghi có kiểm soát} (db);
\node[align=center,font=\small,below=0.45cm of db] {Payments gọi Mock Payment Gateway;\\
OrderExpiryJob kiểm tra đơn quá hạn theo lịch.};
\end{tikzpicture}
\caption{Kiến trúc tổng thể của BrewLite}
\label{fig:architecture}
\end{figure}

\section{Phân chia trách nhiệm backend}
Luồng phân lớp phổ biến là Controller $\rightarrow$ Service
$\rightarrow$ Repository $\rightarrow$ Prisma $\rightarrow$ PostgreSQL.
Controller tiếp nhận request, áp dụng guard/DTO và trả HTTP response;
Service xử lý nghiệp vụ và điều phối transaction; Repository đóng gói
truy vấn; Mapper chuyển đổi dữ liệu sang DTO.

Đây là định hướng tổ chức được thể hiện trong các module, không phải
khẳng định mọi truy vấn đều đã đi qua Repository: \code{OrdersService}
hiện còn trực tiếp dùng Prisma trong transaction tạo đơn, và
\code{OrdersRepository.markAsPaid} chứa cả tác động coupon/loyalty.
Việc đánh giá kiến trúc cần đối chiếu code thay vì chỉ dựa vào tên thư mục.

\begin{table}[H]
\centering\small
\caption{Trách nhiệm các module và thành phần nghiệp vụ}
\label{tab:modules}
\begin{tabularx}{\textwidth}{@{}L{0.22\textwidth}Y@{}}
\toprule\textbf{Thành phần} & \textbf{Trách nhiệm} \\\midrule
Auth / Users & Tài khoản, kiểm tra khách đang hoạt động, JWT, role và quyền sở hữu. \\
Catalog & Sản phẩm, size, topping hợp lệ, giá, công thức và khả năng phục vụ. \\
Orders & Tạo đơn, giá/snapshot, lịch sử/chi tiết, phối hợp transaction và hết hạn. \\
Inventory & Cộng gộp nhu cầu, giữ nguyên liệu, consume/release và kiểm soát version. \\
Payments & Payment attempt, request hash/key, admission, gateway và kết quả thanh toán. \\
Promotions & Thẩm định/khóa coupon và vòng đời giữ/sử dụng/giải phóng lượt. \\
Loyalty & Ledger và balance hiện được xử lý trong OrdersRepository; chưa có module độc lập. \\
\bottomrule
\end{tabularx}
\end{table}

\section{Thiết kế cơ sở dữ liệu}
Schema hiện tại có 18 model; nguồn đối chiếu là
\sourcefile{apps/api/prisma/schema.prisma} và ERD cập nhật trong
\sourcefile{docs/architecture/erd.md}~\cite{schema,erd}. Các bảng được
chia thành các nhóm sau.
{\small\setstretch{1.1}
\begin{longtable}{@{}L{0.20\textwidth}L{0.36\textwidth}L{0.37\textwidth}@{}}
\caption{Các nhóm thực thể trong schema hiện tại}\label{tab:entities}\\
\toprule\textbf{Nhóm} & \textbf{Model} & \textbf{Dữ liệu và quan hệ chính} \\\midrule\endfirsthead
\toprule\textbf{Nhóm} & \textbf{Model} & \textbf{Dữ liệu và quan hệ chính} \\\midrule\endhead
\bottomrule\endfoot
Tài khoản & User & Username duy nhất, password hash, role CUSTOMER/BARISTA, trạng thái hoạt động và loyaltyBalance. Một user có nhiều order. \\
Danh mục & Product, ProductVariant, Topping, ProductTopping & Product có các size; bảng liên kết quy định topping được phép chọn cho từng sản phẩm; đơn giá dùng số nguyên. \\
Nguyên liệu và công thức & Ingredient, RecipeItem, ToppingRecipeItem & Ingredient có đơn vị G/ML/PCS, lượng tồn và version; công thức nối nguyên liệu với size hoặc topping. \\
Đơn và snapshot & Order, OrderItem, OrderItemTopping & Order lưu chủ đơn, tổng tiền, trạng thái và TTL; dòng món/topping giữ tên, giá và tùy chọn tại thời điểm đặt. \\
Giữ chỗ & InventoryReservation, InventoryReservationItem & Reservation gắn order/variant; chi tiết ghi đúng nguyên liệu và lượng đã trừ để hoàn kho. \\
Thanh toán và lịch sử & Payment, OrderStatusHistory & Payment có key, hash, phương thức và kết quả; history lưu trạng thái trước/sau, người đổi và thời gian. \\
Khuyến mãi & Coupon, CouponRedemption & Coupon có loại giảm, thời hạn và giới hạn; redemption gắn user/order và vòng đời giữ lượt. \\
Điểm thưởng & LoyaltyTransaction & Ghi loại giao dịch, số điểm, số dư sau và liên kết order; dùng cùng balance trên User. \\
\end{longtable}}

\subsection{Kho nguyên liệu và công thức}
Không lưu stock của thành phẩm ở ProductVariant. Với một biến thể có tập
nguyên liệu $I$, số ly khả dụng được tính khi đọc:
\[
    \text{availableServings} =
    \min_{i\in I}\left\lfloor\frac{\text{stockQuantity}_i}
    {\text{recipeQuantity}_i}\right\rfloor.
\]
Biến thể chưa có công thức hoặc sử dụng nguyên liệu ngừng hoạt động được
coi là không thể phục vụ. Ví dụ, có 100 g cà phê và công thức cần 10 g/ly
thì giới hạn từ nguyên liệu này là 10 ly; nếu sữa chỉ đủ 6 ly thì giới hạn
chung là 6.

Topping có loại \code{PREPACKAGED} hoặc \code{PREPARED}, đều được biểu
diễn bằng công thức tiêu hao nguyên liệu. Nhu cầu topping được cộng với
nhu cầu đồ uống khi tạo đơn. Khả dụng từng size trên menu chỉ là thông
tin tại thời điểm đọc; quyết định nhận cả giỏ phải kiểm tra lại trong
transaction vì các món có thể dùng chung nguyên liệu.

\subsection{Giữ, sử dụng và hoàn nguyên liệu}
Khi tạo đơn, hệ thống trừ lượng nguyên liệu được giữ và ghi chi tiết vào
reservation trong cùng transaction. Thành công thanh toán chuyển
reservation sang CONSUMED; không trừ nguyên liệu lần thứ hai. Hủy hoặc
hết hạn chỉ hoàn các reservation còn hiệu lực theo đúng lượng đã ghi.

Cập nhật nguyên liệu kiểm tra đồng thời version và lượng tồn đủ. Lượng
tồn không âm được bảo vệ thêm bằng CHECK trong database. Việc tạo đơn sử
dụng mức cô lập Serializable và thử lại tối đa ba lần đối với một số
xung đột được nhận diện.

\subsection{Snapshot và ràng buộc dữ liệu}
\begin{itemize}
    \item Username, SKU, mã coupon và mã đơn có ràng buộc duy nhất;
    quan hệ bảng được bảo vệ bằng khóa ngoại.
    \item Size thuộc một sản phẩm là duy nhất; lượng công thức và
    reservation item phải dương; lượng tồn và giá không âm.
    \item Payment key duy nhất; partial unique index bảo vệ payment
    thành công/đang xử lý theo đơn và giao dịch EARN theo đơn.
    \item Tên món, size, topping và đơn giá được chụp vào đơn; đường/đá
    nullable có miền 0--100, ghi chú tối đa 500 ký tự.
    \item Tiền lưu kiểu số nguyên lớn theo VND và chuyển sang JSON có
    kiểm tra miền an toàn; thời gian lưu \code{timestamptz} theo UTC.
\end{itemize}

\section{Máy trạng thái đơn hàng}
\begin{figure}[H]
\centering
\begin{tikzpicture}[
    state/.style={draw,rounded corners=2pt,minimum width=2.35cm,
        minimum height=0.85cm,font=\scriptsize\ttfamily,align=center},
    real/.style={-{Latex[length=2mm]},thick},
    planned/.style={-{Latex[length=2mm]},thick,dashed}]
\node[state] (pending) at (0,0) {PENDING};
\node[state] (paid) at (3,0) {PAID};
\node[state,dashed] (prep) at (6,0) {PREPARING};
\node[state,dashed] (ready) at (9,0) {READY};
\node[state,dashed] (complete) at (12,0) {COMPLETED};
\node[state] (failed) at (0,-2) {PAYMENT\_FAILED};
\node[state] (cancel) at (6,-2) {CANCELLED};
\draw[real] (pending) -- (paid);
\draw[planned] (paid) -- (prep);
\draw[planned] (prep) -- (ready);
\draw[planned] (ready) -- (complete);
\draw[real] (pending) -- node[right,font=\scriptsize] {Thất bại} (failed);
\draw[real] (failed.west) to[bend left=55]
    node[left,font=\scriptsize] {Thử lại} (pending.west);
\draw[real] (pending.south east) -- node[above,sloped,font=\scriptsize] {Hết hạn} (cancel.north west);
\draw[real] (failed) -- node[below,font=\scriptsize] {Hết hạn} (cancel);
\draw[planned] (paid) -- node[above,sloped,font=\scriptsize] {Khách hủy} (cancel);
\end{tikzpicture}
\caption{Máy trạng thái: nét liền là nhánh đã có; nét đứt là phần cần hoàn thiện}
\label{fig:states}
\end{figure}

Các trạng thái được khai báo trong enum không tự tạo ra hành vi chuyển
trạng thái. Hiện luồng thanh toán, retry và tự hủy do hết hạn đã có code;
nghiệp vụ pha chế/giao món và API khách hủy chưa triển khai đầy đủ.
Thiết kế yêu cầu mọi bước phải có history. Việc thiếu mốc PAID trong
\code{markAsPaid} là khoảng trống hiện trạng, được ghi ở
Chương~\ref{chap:kiem-thu}.

\section{Thiết kế API hiện có}
Tất cả endpoint dưới đây có tiền tố \code{/api/v1}. Bảng chỉ liệt kê
route đã có controller; không tự bổ sung route hủy hoặc route Barista
chưa tồn tại.
{\small\setstretch{1.1}
\begin{longtable}{@{}L{0.11\textwidth}L{0.32\textwidth}L{0.17\textwidth}L{0.31\textwidth}@{}}
\caption{REST API chính đối chiếu từ controller}\label{tab:api}\\
\toprule\textbf{Method} & \textbf{Endpoint} & \textbf{Quyền} & \textbf{Chức năng} \\\midrule\endfirsthead
\toprule\textbf{Method} & \textbf{Endpoint} & \textbf{Quyền} & \textbf{Chức năng} \\\midrule\endhead
\bottomrule\endfoot
GET & \code{/products} & Công khai & Menu và biến thể đang phục vụ. \\
GET & \code{/products/:id} & Công khai & Chi tiết và các topping hợp lệ. \\
POST & \code{/auth/register} & Công khai & Tạo tài khoản khách hàng. \\
POST & \code{/auth/login} & Công khai & Kiểm tra tài khoản và cấp JWT. \\
GET & \code{/auth/me} & JWT & Đọc hồ sơ tài khoản hiện tại. \\
POST & \code{/orders} & CUSTOMER & Tạo đơn, tính giá và giữ chỗ; thành công trả 201. \\
GET & \code{/orders/me} & CUSTOMER & Lịch sử theo tài khoản; limit, cursor và bộ lọc status. \\
GET & \code{/orders/:id} & Chủ đơn & Chi tiết, payment và status timeline. \\
POST & \code{/payments} & CUSTOMER, chủ đơn & Thanh toán; bắt buộc header Idempotency-Key. \\
POST & \code{/coupons/validate} & CUSTOMER & Kiểm tra coupon và trả báo giá. \\
GET & \code{/health} & Công khai & Kiểm tra API và khả năng truy cập database. \\
\end{longtable}}

\section{Thiết kế các màn hình chính}
Luồng giao diện: \code{/menu} $\rightarrow$ modal chi tiết
$\rightarrow$ CartDrawer $\rightarrow$ \code{/checkout}
$\rightarrow$ \code{/orders/confirmation}. Các route \code{/login} và
\code{/register} xử lý tài khoản; \code{/orders} và
\code{/orders/[id]} cung cấp lịch sử/chi tiết.

Giỏ giữ trạng thái cục bộ để người dùng chưa đăng nhập vẫn chọn món.
Khi checkout, phiên đăng nhập cung cấp token; giao diện gửi ID và tùy
chọn thay vì tự quyết định tiền phải trả. Sau thất bại thanh toán, giỏ
vẫn giữ nguyên; sau thành công, khách chuyển đến xác nhận theo mã đơn
server đã tạo.
